\documentclass[11pt]{article}
\usepackage[T1]{fontenc}
\usepackage[utf8]{inputenc}
\usepackage{lmodern}
\usepackage[letterpaper,margin=0.85in]{geometry}
\usepackage{microtype}
\usepackage{booktabs}
\usepackage{longtable}
\usepackage{enumitem}
\usepackage{xcolor}
\usepackage[hidelinks]{hyperref}
\setlength{\parindent}{0pt}
\setlength{\parskip}{5pt}
\setlist{itemsep=2pt,topsep=3pt}
\newcommand{\code}[1]{\texttt{\detokenize{#1}}}
\title{Consensus Response to the Patched Master Collaborative Reference}
\author{Review prepared for Edward Wang's Kitware internship record}
\date{July 15, 2026}
\begin{document}
\maketitle

\section*{Overall assessment}
The patch is careful, constructive, and substantively correct. I accept the core
scientific framing and almost all factual corrections. The remaining disagreements
concern evidence labels and publication acceptance, not the paper's thesis.

\begin{quote}\itshape
EEE detects reproducibility gaps. The follow-on work treats benchmark reproduction
as layered system identification: it attempts to attribute and sometimes close
disagreement by reconstructing recipe, deployment, execution-instrument, and artifact
provenance, while recognizing that some historical measurements are non-identifiable
from the surviving record examined.
\end{quote}

The near-term manuscript should remain a hybrid position/systems-and-measurement paper.
OLMo is the partially recoverable modern case; classic HELM supplies proxy-era and
non-identifiable historical cases; a narrow Qwen study is optional rather than a
dependency.

\section{Decisions on the patch}

\subsection*{A1. Commit \code{86ec84af}}
\textbf{Accept the factual correction; modify the evidence label.} The direct live
repository observation is accepted, but the package does not include the Git object,
a Git bundle, or an immutable remote reference. The master therefore uses
\emph{collaborator-verified, artifact not packaged}, rather than \emph{Established}.

\subsection*{A2. Result stores}
\textbf{GPT-OSS: accept report freshness; reject preservation-only as the remaining
step.} A report timestamp does not establish run-artifact freshness, comparison
freshness, or a complete provenance manifest. The cells are promising candidate
results pending acceptance and preservation.

\textbf{OLMo: accept.} The reported local-halving signature makes the store stale
relative to corrected comparison logic. Regenerate before citation.

\textbf{RedPajama: accept with scope qualification.} The live inspection strengthens
the both-era proxy-execution claim. It does not prove either proxy is the original
instrument. Copy, hash, and preserve the complete inputs and provenance.

\textbf{Qwen: accept.} No result store is present; the study remains pending and is
not a near-term paper dependency.

\subsection*{A3. Probe scope and the 59\% estimate}
\textbf{Accept the IFEval-only deployment-match scope}, with a requirement to package
sweep manifests and outputs. \textbf{Accept the 59\% value only as an internal
planning estimate.} Regenerate it from a frozen manifest with an explicit numerator,
denominator, eligibility rule, and date before publication.

\subsection*{B1. Dense OLMo-2 HF divergence}
\textbf{Accept.} Both genuine execution-instrument divergence and a dense-model probe
artifact remain live explanations. The evidence locates a divergence, not its unique
root cause.

\subsection*{B2. \code{HuggingFaceClient} wording}
\textbf{Accept with a source-version qualification.} ``Recorded as a
\code{HuggingFaceClient} deployment'' is precise; it must not imply a uniquely known
historical package build.

\subsection*{B3. Drift ratio}
\textbf{Accept.} The correct description is approximately 25--35x (about 30x), not two
orders of magnitude.

\subsection*{B4. Canonical ledger}
\textbf{Accept and implement.} The source of truth is
\code{chronicle/data/master_claim_evidence_ledger_2026-07-15c.csv}. Prose tables are
snapshots. Store status is separately represented in
\code{chronicle/data/store_status_ledger_2026-07-15c.csv}.

\subsection*{B5. Vendor-named front matter}
\textbf{Accept in principle.} The front summary now refers to independent
reconstructions. A concise evidence-status chapter remains near the front because it
governs the chronology; detailed correspondence is preserved under \code{validation/}.

\subsection*{B6. Packaging and compilation}
\textbf{Accept and implement fully.} The structured package contains both commit-ledger
appendices, an accurate README, patch materials, canonical ledgers, and compiled
primary documents.

\section{Additional concerns from the independent review}

\subsection*{Report freshness is not result freshness}
Record separately: (1) run-artifact freshness, (2) comparison freshness, and (3)
report freshness. A recent report timestamp establishes only the third unless a
manifest links all three.

\subsection*{Add an orthogonal reproduction-target axis}
The six-category taxonomy explains \emph{why} executions disagree. The paper should
also distinguish: artifact reconstruction, procedural reproduction, and claim-level
robustness. Official released artifacts remain ground truth for the historical score;
reruns test procedure stability and identifiability.

\subsection*{Scope non-identifiability}
Use ``non-identifiable from the surviving public and repository evidence examined,''
not an absolute claim that no private archive could resolve the instrument.

\subsection*{Keep EvalAudit's claim at the capability level}
The current defensible claim is support for exact-spec replay, provenance preservation,
layered comparison, and controlled diagnosis. Do not claim it guarantees reproducible
evaluation or ``fixes HELM'' without system validation.

\subsection*{Preserve live-only evidence through attestations}
For each live-store or post-archive Git claim, preserve identifiers and hashes, times,
code and dirty state, run specs and targets, model/deployment revisions, environment
and hardware identity, and generation/comparison commands.

\section{Final consensus}
There is no remaining substantive disagreement about the paper's core thesis or
near-term story.

\begin{quote}\bfseries
A public LLM benchmark recipe is not a complete experimental record. Benchmark
reproduction is a layered system-identification problem spanning recipe, deployment,
execution instrument, and artifact interpretation. Controlled reconstruction can
attribute and sometimes close apparent gaps, while some historical measurements
remain non-identifiable from the surviving evidence examined.
\end{quote}

The narrative is: EEE detects; EvalAudit replays, preserves, compares, and diagnoses;
OLMo is partially recoverable pending ordinary-path confirmation; RedPajama compares
declared proxy eras; GPT-J/GPT-NeoX/OPT illustrate origin-instrument
non-identifiability; the paper derives a minimum provenance standard; and Qwen is an
optional current-model demonstration.

The remaining work is operational: regenerate OLMo, preserve and audit GPT-OSS and
RedPajama, package sweep manifests, complete or document failure of OLMo confirmation,
regenerate denominators, and keep the canonical ledger synchronized.

With those qualifications incorporated, the chronology and paper strategy have
reached consensus.
\end{document}
